\section{Additional experimental details}
\label{app:experimental_details}

This appendix records the paper-facing experimental protocol. The main experiments ask whether sparse Koopman autoencoders learn useful latent supports without basin supervision. The model is never given basin labels, basin counts, or trajectory-to-basin assignments when training the encoder, decoder, latent transition, support rules, or periodic rollouts. Benchmark basin metadata is used only after training to define evaluation slices, construct controlled diagnostic interventions, and score support--basin agreement. The only training-time exception is architectural: the block-diagonal and soft-block transition diagnostics use benchmark metadata to set the fixed number of transition groups, but they still do not receive trajectory labels or support labels.

\paragraph{Benchmarks.}
The controlled benchmark contains \(15\) two-dimensional multibasin systems, listed with their equations and attractor metadata in \cref{app:fixed17_inventory}. These systems are integrated to produce stored observation sequences; the learner observes only stored states, not vector fields, integration substeps, basin labels, or basin assignments. The external forecasting stress test uses the \(10\) three-dimensional Dysts \(dt{\times}30\) systems listed in \cref{app:dysts_inventory_full}. For Dysts, stored observations are generated at \(30\) times each system's native integration interval and coordinates are standardized after trajectory generation.

\paragraph{Model rows.}
The six model rows in \cref{tab:forecasting_main,tab:fixed17_alignment} are LISTA, LISTA--BD, LISTA--SB, Sparse MLP, Sparse MLP--BD, and Dense MLP. All use latent dimension \(d_z=256\). LISTA rows use shrinkage to produce exact zeros; MLP sparse rows use the sparsity penalty in \cref{eq:loss_latent_terms}; Dense MLP removes the intended sparsity mechanism and sets the sparsity coefficient to zero. The BD rows use a block-diagonal latent transition. The SB row keeps a dense transition but adds an off-block penalty. All rows use a linear decoder dictionary as the state decoder. The controlled multibasin LISTA rows use threshold scale \(\alpha=0.15\), two shrinkage-refinement loops, and a sign-split final operation. The Dysts \(dt{\times}30\) LISTA rows use the same threshold scale with one refinement loop and a ReLU final operation. MLP encoders use two hidden layers of width \(64\); Dense MLP uses a tanh encoder without the sparse output gate.

\paragraph{Controlled multibasin training.}
For every controlled system, each model row is trained for seeds \(0,\ldots,14\). Training uses rollout length \(L=8\) windows, meaning \(9\) adjacent stored states per window, for \(200{,}000\) optimization steps with minibatches of \(256\) windows. The optimizer is AdamW. Encoder and decoder parameters use learning rate \(5{\times}10^{-5}\), the latent transition uses learning rate \(5{\times}10^{-6}\), and non-transition parameters use weight decay \(10^{-4}\). In \cref{eq:training_objective}, the loss weights are \(\lambda_{\rm pred}=1\), \(\lambda_{\rm rec}=0.03\), \(\lambda_{\rm lin}=1\), and \(\lambda_{\rm sp}=3{\times}10^{-3}\) for sparse rows; Dense MLP uses \(\lambda_{\rm sp}=0\). Soft-block rows add the off-block \(L_1\) transition penalty with weight \(10^{-4}\).

All controlled rows use the same label-free boundary-emphasized reset distribution. Before training, \(4096\) candidate initial states are probed for \(32\) steps. Candidates are scored by short-horizon perturbation sensitivity and residual late-rollout motion, and the top \(1024\) states are retained. The scoring probe uses four perturbations, perturbation scale \(0.04\), a late-transient window of \(8\) steps, and transient weight \(0.5\). Half of training windows start from this retained pool with jitter scale \(0.25\). The retained pool is intended to expose boundary-adjacent and transient regions that would otherwise be rare under ordinary resets; it is built without basin labels.

\paragraph{Dysts \(dt{\times}30\) training.}
Dysts rows are trained for the same seeds, \(0,\ldots,14\), with latent dimension \(d_z=256\), minibatches of \(256\), rollout length \(L=10\) windows, and \(100{,}000\) optimization steps. Dysts trajectories come from deterministic native caches with \(200\) cached trajectories, \(30{,}000\) stored observations per trajectory, and a \(2{,}000\)-step warm-up. The first cached trajectory starts from the Dysts default initial condition; remaining trajectories perturb that state with Gaussian noise at scale \(0.2\) times the coordinate-wise Dysts standard deviation. Train, validation, and test caches use separate deterministic namespaces, so held-out test windows are disjoint from training windows and reusable across model seeds. LISTA rows use the controlled-benchmark learning rates. MLP rows use learning rates \(10^{-4}\) for encoder/decoder parameters and \(10^{-5}\) for the latent transition. Sparse Dysts rows use \(\lambda_{\rm sp}=6{\times}10^{-3}\), while Dense MLP uses \(\lambda_{\rm sp}=0\).

\paragraph{Checkpoint selection.}
Checkpoint selection is fixed before test evaluation. Every \(500\) optimization steps, and again at the final step, the current model is rolled out for \(200\) stored steps from \(16\) fixed held-out initial states using every-step re-encoding. The checkpoint with the lowest final validation error under this proxy is used for reported evaluations. Final tables aggregate completed seeds; they do not select the best seed.

\paragraph{Training recipe selection.}
Exploratory sweeps varied LISTA depth, shrinkage scale, sparsity coefficient, sign-split encodings, transition structure, and MLP sparsity controls. These sweeps used shorter budgets or broader system lists and are not pooled with the final evidence. The paper-facing protocol is the fixed recipe in this appendix: \(15\) seeds per retained system, the training budgets above, held-out evaluation splits, and the statistical units in \cref{app:statistical_testing}.

\paragraph{Forecasting experiments.}
Forecasting is evaluated on held-out trajectories at the stored observation cadence. A horizon \(h\) always means \(h\) stored observation steps, not a common physical time across systems. The no-reencoding rollout applies the learned latent transition for the full horizon before decoding. The periodic rollout advances the model's own predicted latent state, decodes the predicted state at a fixed period \(m\), and re-encodes that predicted state; it never uses the true future state. Controlled multibasin forecasting uses \(100\) held-out initial conditions for each system and seed and reports the best score over the fixed period grid \(m\in\{10,25,50,100\}\). Dysts long-horizon forecasting uses \(100\) held-out test-cache windows for each system and seed and reports the best score over \(m\in\{10,25,50,100,150,200\}\). These grids are fixed before aggregation and are not tuned separately for each test trajectory.

\paragraph{Support--basin alignment.}
Support alignment is computed after training from held-out controlled-system states. For each state \(x\), the basin-depth margin is the distance to the second-nearest benchmark attractor center minus the distance to the nearest attractor center. Within each represented benchmark basin, the evaluator keeps the top quartile of held-out states by this margin. This per-basin deep slice uses ground-truth geometry only to define the evaluation population. It is the clean slice for testing basin-support alignment because boundary states can legitimately have unstable or mixed supports.

On this slice, the evaluator encodes states and constructs \(S_{\rm abs}(x)=\{i:|z_i|>10^{-3}\}\). Exact Boolean masks are grouped into \(F_{\rm abs}\) support families by the deterministic greedy Jaccard rule in \cref{sec:method_supports,app:support_definitions} with threshold \(0.5\). The main directional metric is \(H(B\mid F_{\rm abs})\), where \(B\) is the withheld benchmark basin label; lower values mean that knowing the learned support family leaves less uncertainty about basin identity. The companion count \(\overline{|F_{\rm abs}|}\) is descriptive and is interpreted relative to the number of represented basins, not as a monotone score to maximize.

\paragraph{Wrong-support functional diagnostic.}
The cross-system wrong-support diagnostic asks whether the selected active coordinates matter for prediction. For each testable controlled system and basin, the evaluator constructs a canonical deep-slice \(S_{\rm abs}\) mask from the most frequent exact mask in that basin. During a held-out rollout, the inferred mask is replaced with a canonical mask from a different basin and held fixed during the masked latent rollout. \Cref{tab:fixed17_alignment} reports wrong-support/base error ratios at \(H=1\) and \(H=20\); larger ratios mean that using a basin-inconsistent support is more damaging.

\paragraph{Support-coordinate intervention figure.}
The representative coordinate-intervention experiment isolates active-coordinate identity from coefficient values. Starting from a held-out basin-interior state, the model encodes \(z_0=\Enc(x_0)\), perturbs only \(z_0\), and then runs the ordinary learned transition and decoder without retraining. In the coordinate-dropping condition, the active entries are ranked by \(|z_{0,i}|\) and the top \(k\in\{1,2,3,5,10\}\) are set to zero before rollout. In the random-support condition, active coefficient values are preserved but reassigned to randomly chosen inactive coordinates. The figure and table use \(100\) held-out basin-interior starts; the random-support condition uses \(20\) random inactive-coordinate reassignments for each start. Errors are accumulated through \(H=21\), with the full horizon curves shown in \cref{fig:support_coordinate_interventions} and the \(H=21\) summary in \cref{tab:support_coordinate_interventions_h21}.

\paragraph{Periodic support-refresh diagnostic.}
The support-refresh experiment asks whether periodic re-encoding can re-select the support family appropriate to a new region of state space after a controlled transfer. The support object is \(F_{\rm top8}\), built by taking the eight largest-magnitude latent coordinates and merging exact top-\(8\) masks with the same Jaccard-family rule. Each valid comparison starts from a controlled source-to-target transfer generated from benchmark annotations. Between scheduled refresh events, the latent state is advanced only with the learned global transition \(K\). At each event, the stale family implied by the propagated latent is compared with the family obtained by encoding the current controlled-transfer state. The reported cells in \cref{tab:support_refresh} are target-family fractions before re-encoding \(\to\) after re-encoding for fixed periods \(m\in\{1,5,10,20\}\), plus a no-transfer false-target control. Basin labels are used to construct and score transfer instances, not to train a support-conditioned rollout.

\paragraph{Post-hoc routing diagnostic.}
The appendix support-conditioned predictor experiment is a second-stage diagnostic, not part of the main training procedure. After a Koopman autoencoder checkpoint is fixed, the encoder, decoder, and original global \(K\) are frozen. Routes are built from \(S_{\rm top8}\) masks and merged into \(F_{\rm top8}\) support families with Jaccard threshold \(0.40\). Local transition maps are fit by support family on a disjoint fitting split of \(256\) held-out trajectories of length \(256\), giving \(256\times256=65{,}536\) adjacent latent transitions per checkpoint and support definition. A family with fewer than \(50\) current-state fitting transitions does not receive a trainable local map. For each retained family \(c\), the centered affine map \(T_c(z)=\bar z_c+K_c(z-\bar z_c)\) is initialized from the checkpoint's global \(K\) and optimized with decoded rollout MSE while the original autoencoder remains frozen. Route-balanced minibatches prevent the most frequent support families from monopolizing the second-stage updates. Evaluation uses a separate set of \(128\) held-out rollouts; at \(H=1000\), this contributes \(128{,}000\) one-step target positions along evaluated rollouts. Routes are recomputed from the current predicted latent state during rollout, and the routed rollout periodically decodes and re-encodes its own predicted state with period \(5\). If no fitted local route is available, the evaluator falls back to the same model's global transition. The reported routed/global ratio therefore compares a support-routed rollout with the same frozen checkpoint's own global rollout.

\paragraph{Aggregation and significance.}
Point estimates in forecasting and support tables summarize seeds within each system by interquartile mean and then average system summaries across systems. Multibasin forecasting, support-alignment, and wrong-support significance annotations use paired seed-level effects within each controlled system against Dense MLP, with one-sided alternatives fixed by the table direction and Holm correction across eligible systems. Dysts forecasting uses each Dysts system as the independent unit after within-system seed-IQM summarization and applies a paired one-sided exact sign test with Holm correction. Support-refresh and coordinate-intervention panels are descriptive mechanism diagnostics. Full statistical definitions and denominator conventions are in \cref{app:statistical_testing}.


\paragraph{Hardware and resource details.}
Training was done on a SLURM cluster; RTX8000 was the main GPU being used.
Training jobs load CUDA \(12.6.0\) and request one cluster GPU, \(4\) CPU cores,
and \(16\) GB host memory per active training allocation.  Controlled
multibasin training uses single-run GPU allocations with a \(24\)-hour time
limit.  Dysts \(dt{\times}30\) training uses packed GPU allocations with up to
\(12\) runs per allocation, one GPU, \(4\) CPU cores, \(16\) GB memory, and a
\(72\)-hour time limit per packed allocation.  Dysts cache construction uses
CPU allocations with \(4\) CPU cores, \(16\) GB memory, and a \(24\)-hour time
limit.  Dysts long-horizon evaluation uses CPU allocations with \(4\) CPU
cores, \(24\) GB memory, and a \(6\)-hour time limit for the paper-facing
\(dt{\times}30\) evaluation queue.  Controlled support diagnostics and
support-routed analyses use CPU allocations with \(4\) CPU cores, \(16\)--\(24\)
GB memory, and \(8\)--\(12\)-hour time limits, depending on the diagnostic.